% =========================================================
% Limits Superior and Inferior of Functions — Canonical Notes
% Chapter: continuity
% Volume III · Analysis
% =========================================================

\section{Limits Superior and Inferior of Functions}
\label{sec:limsup-liminf-functions}

% ---------------------------------------------------------
% Toolkit
% ---------------------------------------------------------

\begin{tcolorbox}[title={Toolkit — Limits Superior and Inferior of Functions}]
\begin{itemize}
  \item $\limsup_{x \to x_0} f(x) = \inf_{\delta>0} \sup\{f(x) : 0 < |x-x_0| < \delta,\; x \in E\}$.
  \item $\liminf_{x \to x_0} f(x) = \sup_{\delta>0} \inf\{f(x) : 0 < |x-x_0| < \delta,\; x \in E\}$.
  \item Both always exist in $[-\infty, +\infty]$.
  \item $\limsup \geq \liminf$ always.
  \item $\lim_{x \to x_0} f(x)$ exists $\iff$ $\limsup = \liminf \in \mathbb{R}$.
  \item Direct analogue of the sequence versions: replace tails with punctured neighborhoods.
\end{itemize}
\end{tcolorbox}

% ---------------------------------------------------------
% Definition — Lim Sup and Lim Inf of a Function
% ---------------------------------------------------------

\begin{tcolorbox}[title={Definition — Limits Superior and Inferior of a Function}]
\begin{definition}[Limits Superior and Inferior of a Function]
\label{def:limsup-liminf-function}
Let $f : E \to \mathbb{R}$ and let $x_0$ be a point of accumulation
of $E$. Define
\[
  \limsup_{x \to x_0} f(x)
  \;:=\;
  \inf_{\delta > 0}
  \sup\bigl\{\,f(x) : x \in E,\; 0 < |x - x_0| < \delta\,\bigr\}
\]
and
\[
  \liminf_{x \to x_0} f(x)
  \;:=\;
  \sup_{\delta > 0}
  \inf\bigl\{\,f(x) : x \in E,\; 0 < |x - x_0| < \delta\,\bigr\}.
\]
\end{definition}
\end{tcolorbox}

\begin{remark*}[Standard quantified statement]
\[
  \limsup_{x \to x_0} f(x)
  = \inf_{\delta > 0}
    \sup\{f(x) : x \in E \cap (V_\delta(x_0) \setminus \{x_0\})\}.
\]
\end{remark*}

\begin{remark*}[Why the inf of sup]
As $\delta$ decreases, the set over which the sup is taken shrinks,
so the sup is nonincreasing in $\delta$. The $\limsup$ is the
infimum of this nonincreasing function as $\delta \to 0^+$, i.e.,
the limit. The $\liminf$ is dual: the sup of a nondecreasing
infimum.
\end{remark*}

\begin{remark*}[Connection to sequence limsup and liminf]
The sequence version: $\limsup_{n \to \infty} a_n =
\inf_N \sup_{n \geq N} a_n$. The function version is identical
in spirit: replace the tail $\{a_n : n \geq N\}$ with the
punctured neighborhood $\{f(x) : 0 < |x - x_0| < \delta\}$.
The analogy is exact.
\end{remark*}

\begin{remark*}[Interpretation]
Both $\limsup$ and $\liminf$ always exist in $[-\infty, +\infty]$.
They measure, respectively, the largest and smallest cluster values
of $f(x)$ as $x \to x_0$. If $\limsup = \liminf$, the function
has a well-defined limit; if they differ, the function oscillates
between distinct accumulation values. The canonical example is
$f(x) = \sin(1/x)$ at $x_0 = 0$: $\limsup = 1$ and $\liminf = -1$.
\end{remark*}

% ---------------------------------------------------------
% Proposition — Lim Sup Is At Least Lim Inf
% ---------------------------------------------------------

\begin{proposition}[Lim Sup Is At Least Lim Inf]
\label{prop:limsup-geq-liminf-function}
For any $f : E \to \mathbb{R}$ and any accumulation point $x_0$ of
$E$,
\[
  \limsup_{x \to x_0} f(x) \;\geq\; \liminf_{x \to x_0} f(x).
\]
\end{proposition>

\begin{remark*}[Standard quantified statement]
\[
  \forall f : E \to \mathbb{R},\;
  \forall x_0 \text{ accumulation point of } E :\;
  \limsup_{x \to x_0} f(x) \geq \liminf_{x \to x_0} f(x).
\]
\end{remark*>

\begin{remark*}[Interpretation]
For any fixed $\delta > 0$, $\sup \geq \inf$ on the punctured
neighborhood. Taking the infimum over $\delta$ on the left and
the supremum over $\delta$ on the right preserves the inequality.
\end{remark*}

% ---------------------------------------------------------
% Proposition — Limit Exists iff Lim Sup Equals Lim Inf
% ---------------------------------------------------------

\begin{proposition}[Limit Exists iff Lim Sup Equals Lim Inf]
\label{prop:limsup-liminf-limit-criterion}
Let $f : E \to \mathbb{R}$ and let $x_0$ be an accumulation point
of $E$. Then $\lim_{x \to x_0} f(x)$ exists as a finite real number
if and only if
\[
  \limsup_{x \to x_0} f(x) = \liminf_{x \to x_0} f(x) \in \mathbb{R}.
\]
In that case, all three quantities are equal.
\end{proposition>

\begin{remark*}[Standard quantified statement]
\[
  \lim_{x \to x_0} f(x) = L \in \mathbb{R}
  \;\iff\;
  \limsup_{x \to x_0} f(x) = L = \liminf_{x \to x_0} f(x).
\]
\end{remark*>

\begin{remark*}[Interpretation]
The exact analogue of the sequence criterion: $\lim_{n \to \infty}
a_n$ exists iff $\limsup a_n = \liminf a_n$. The limit exists when
the upper and lower envelopes collapse to the same finite value.
If they disagree — as with $\sin(1/x)$ at $0$ where $\limsup = 1
\neq -1 = \liminf$ — no limit exists.
\end{remark*>
